\subsection{Paired calendar-block score uncertainty}
\label{sec:block_score_uncertainty}
For each test fold, contiguous 90-day calendar blocks retain both the event
log-intensity contrast and the corresponding integrated-rate contrast,
including empty blocks and terminal exposure. We resample these paired block
scores 49,999 times with replacement to obtain percentile 95\% intervals for
the total contrast. A two-sided paired sign-flip reference independently
reverses each block contrast. Training estimates remain fixed in these
conditional comparisons. Pooled totals resample separately within each disjoint
fold before summing. We also repeat the calculation with 180- and 365-day blocks.

Calendar splitting does not make these blocks independent or strictly
exchangeable. In the New Zealand model, 61 of the 195 finite-mixture modes
decay by less than $1\%$ over $24{,}394$ days, with minimum rate
$e^{-38.8}\,\mathrm{day}^{-1}$; their aggregate weight is small, but the
state is carried across block boundaries. Triggered sequence clustering and
fitted nonstationary rates provide further dependence across months and years.
An independently reversed sign for each observed block is therefore a
reference calculation under an additional symmetry/exchangeability assumption,
not a justified randomization of independent earthquake sequences.

Small nominal sign-flip values and wide block intervals address different
features of the same contrasts. For example, Fold~1 has nominal $p=0.00104$
and interval $[-25.526,-7.459]$ nats, while its strongest subinterval loss is
concentrated in 1995--1998. The sign-flip reference counts directional
extremeness under independent reversals; the bootstrap expresses magnitude
variability under its chosen block resampling. Neither operation establishes
long-memory calibration. Prioritize the agreement or disagreement of the
90-, 180- and 365-day interval summaries and the fold-level directions over a
nominal 90-day $p$-value. The intervals also remain conditional sensitivity
summaries: their nominal coverage is not proved under the catalogue's full
dependence and fitted-parameter uncertainty. Section~\ref{sec:diagnostic_scope}
states the common inferential scope.

\begin{table}[htbp]
\centering
\begin{tabular}{lrrr}
\toprule
Fold & Contrast (nats) & Paired 95\% interval & Nominal sign-flip $p$ \\
\midrule
1 & -16.518 & $[-25.526,-7.459]$ & 0.00104 \\
2 & -0.075 & $[-1.320,1.215]$ & 0.91248 \\
3 & -1.030 & $[-2.087,0.004]$ & 0.05766 \\
\bottomrule
\end{tabular}
\caption{Shared inhibition minus the uninhibited Mittag--Leffler model: paired calendar-block uncertainty. Negative contrasts favour the uninhibited model.}
\label{tab:nz_score_uncertainty}
\end{table}

Fold~1 contributes 93.7\% of the total deficit and its intervals exclude zero at all three block lengths. Fold~2's intervals contain zero throughout: its 180- and 365-day intervals are $[-1.109,1.014]$ and $[-1.511,1.324]$. Fold~3 is sensitive to the dependence scale: the 90-day interval includes zero, whereas the 180- and 365-day intervals are $[-2.185,-0.091]$ and $[-2.149,-0.225]$; the annual-block nominal sign-flip value is $p=0.015$. This change in zero inclusion is part of the block-size sensitivity, not a validated long-memory rejection.
